% !TEX program = xelatex
\documentclass[a4paper,12pt]{article}

\usepackage{polyglossia}
\setmainlanguage{russian}
\usepackage{fontspec}
% Если Times New Roman недоступен, замените на: TeX Gyre Termes
\setmainfont{Times New Roman}
\usepackage[margin=1.5cm]{geometry}

\usepackage{graphicx}
\usepackage{tikz}
\usetikzlibrary{calc}

\begin{document}
	\centering
	
	% ===== НАСТРОЙКИ =====
	\def\FrontImage{front.jpg}   % <-- ИМЯ ВАШЕГО ФОТО
	\def\Cols{6}\def\Rows{4}     % сетка 6x4
	
	% Зона сетки в долях ширины/высоты изображения (0..1):
	\def\XL{0.05}\def\XR{0.50}   % левая/правая границы по ширине
	\def\YT{0.18}\def\YB{0.78}   % верх/низ по высоте
	
	% Уровни по высоте (мм) — 100, 220, 340, 460 (формула: 100 + 120*r, r=0..3)
	
	% ---------------- КАРТИНКА 1: КЛИН ----------------
	\begin{tikzpicture}
		\node[inner sep=0] (im) {\includegraphics[width=\linewidth]{\FrontImage}};
		% Углы изображения:
		\path (im.south west) coordinate (SW);
		\path (im.north east) coordinate (NE);
		
		% Границы зоны (через доли 0..1):
		\coordinate (XLpt) at ($(SW)!\XL!(NE)$);
		\coordinate (XRpt) at ($(SW)!\XR!(NE)$);
		\coordinate (YTpt) at ($(SW)!\YT!(NE)$);
		\coordinate (YBpt) at ($(SW)!\YB!(NE)$);
		
		% Прямоугольник зоны:
		\coordinate (GLL) at (XLpt |- YBpt);  % нижний левый
		\coordinate (GUR) at (XRpt |- YTpt);  % верхний правый
		
		% Рамка и сетка:
		\draw[yellow, line width=1.8pt] (GLL) rectangle (GUR);
		\foreach \c in {1,...,\numexpr\Cols-1\relax}{
			\coordinate (V\c) at ($(GLL)!\c/\Cols!(GUR -| GLL)$);
			\draw[white, opacity=.85, line width=.8pt] (V\c) -- (V\c |- GUR);
		}
		\foreach \r in {1,...,\numexpr\Rows-1\relax}{
			\coordinate (H\r) at ($(GLL)!\r/\Rows!(GUR |- GLL)$);
			\draw[white, opacity=.85, line width=.8pt] (H\r) -- (H\r -| GUR);
		}
		
		% Заголовок:
		\node[anchor=south west, fill=black!60, text=white, inner sep=3pt]
		at ($(GUR)+(0,-10mm)$) {CRASH3 матрица 6×4 (мм) — клин};
		
		% Ячейки и значения (клин: scale = max(0, 1 - (c+0.5)/Cols) ):
		\foreach \r in {0,...,3}{% уровень по высоте
			\pgfmathsetmacro{\Lmm}{100 + 120*\r}  % 100,220,340,460
			\pgfmathsetmacro{\yTop}{(\r+1)/\Rows}
			\pgfmathsetmacro{\yBot}{\r/\Rows}
			\coordinate (Yt) at ($(GLL)!\yTop!(GUR |- GLL)$);
			\coordinate (Yb) at ($(GLL)!\yBot!(GUR |- GLL)$);
			\foreach \c in {0,...,5}{% сечение по ширине
				\pgfmathsetmacro{\xL}{\c/\Cols}
				\pgfmathsetmacro{\xC}{(\c+0.5)/\Cols}
				\pgfmathsetmacro{\xR}{(\c+1)/\Cols}
				\coordinate (X0) at ($(GLL)!\xL!(GUR -| GLL)$);
				\coordinate (Xc) at ($(GLL)!\xC!(GUR -| GLL)$);
				\coordinate (X1) at ($(GLL)!\xR!(GUR -| GLL)$);
				
				\pgfmathsetmacro{\Scale}{max(0, 1 - (\c+0.5)/\Cols)}
				\pgfmathsetmacro{\V}{round(\Lmm * \Scale)}    % мм
				\pgfmathsetmacro{\Perc}{min(100, max(0, round(100*\V/460)))}
				
				% Заливка и контур:
				\fill[red!\Perc!black, opacity=.35] (X0|-Yb) rectangle (X1|-Yt);
				\draw[white, line width=.8pt, opacity=.95] (X0|-Yb) rectangle (X1|-Yt);
				
				% Подпись значения:
				\node[fill=black!55, text=white, inner sep=2pt, rounded corners=1pt]
				at (Xc|-$(Yb)!0.5!(Yt)$) {\pgfmathprintnumber{\V}};
			}
		}
	\end{tikzpicture}
	
	\clearpage
	
	% ---------------- КАРТИНКА 2: РАВНОМЕРНО ----------------
	\begin{tikzpicture}
		\node[inner sep=0] (im) {\includegraphics[width=\linewidth]{\FrontImage}};
		\path (im.south west) coordinate (SW);
		\path (im.north east) coordinate (NE);
		
		\coordinate (XLpt) at ($(SW)!\XL!(NE)$);
		\coordinate (XRpt) at ($(SW)!\XR!(NE)$);
		\coordinate (YTpt) at ($(SW)!\YT!(NE)$);
		\coordinate (YBpt) at ($(SW)!\YB!(NE)$);
		
		\coordinate (GLL) at (XLpt |- YBpt);
		\coordinate (GUR) at (XRpt |- YTpt);
		
		\draw[yellow, line width=1.8pt] (GLL) rectangle (GUR);
		\foreach \c in {1,...,\numexpr\Cols-1\relax}{
			\coordinate (V\c) at ($(GLL)!\c/\Cols!(GUR -| GLL)$);
			\draw[white, opacity=.85, line width=.8pt] (V\c) -- (V\c |- GUR);
		}
		\foreach \r in {1,...,\numexpr\Rows-1\relax}{
			\coordinate (H\r) at ($(GLL)!\r/\Rows!(GUR |- GLL)$);
			\draw[white, opacity=.85, line width=.8pt] (H\r) -- (H\r -| GUR);
		}
		
		\node[anchor=south west, fill=black!60, text=white, inner sep=3pt]
		at ($(GUR)+(0,-10mm)$) {CRASH3 матрица 6×4 (мм) — равномерно};
		
		% Ячейки и значения (равномерно: scale = 1):
		\foreach \r in {0,...,3}{
			\pgfmathsetmacro{\Lmm}{100 + 120*\r}
			\pgfmathsetmacro{\yTop}{(\r+1)/\Rows}
			\pgfmathsetmacro{\yBot}{\r/\Rows}
			\coordinate (Yt) at ($(GLL)!\yTop!(GUR |- GLL)$);
			\coordinate (Yb) at ($(GLL)!\yBot!(GUR |- GLL)$);
			\foreach \c in {0,...,5}{
				\pgfmathsetmacro{\xL}{\c/\Cols}
				\pgfmathsetmacro{\xC}{(\c+0.5)/\Cols}
				\pgfmathsetmacro{\xR}{(\c+1)/\Cols}
				\coordinate (X0) at ($(GLL)!\xL!(GUR -| GLL)$);
				\coordinate (Xc) at ($(GLL)!\xC!(GUR -| GLL)$);
				\coordinate (X1) at ($(GLL)!\xR!(GUR -| GLL)$);
				
				\pgfmathsetmacro{\V}{round(\Lmm)}    % мм
				\pgfmathsetmacro{\Perc}{min(100, max(0, round(100*\V/460)))}
				
				\fill[red!\Perc!black, opacity=.35] (X0|-Yb) rectangle (X1|-Yt);
				\draw[white, line width=.8pt, opacity=.95] (X0|-Yb) rectangle (X1|-Yt);
				
				\node[fill=black!55, text=white, inner sep=2pt, rounded corners=1pt]
				at (Xc|-$(Yb)!0.5!(Yt)$) {\pgfmathprintnumber{\V}};
			}
		}
	\end{tikzpicture}
	
	\vspace{2mm}
	{\small
		Точное позиционирование: корректируйте доли \texttt{\XL, \XR, \YT, \YB}. 
		Напр., \texttt{\XL=0.10, \XR=0.52, \YT=0.22, \YB=0.78}.
		Чтобы рисовать только один вариант — удалите лишний блок \texttt{tikzpicture}.
	}
	
\end{document}